\cleardoublepage
\chapter*{Historical Note}
\label{nota-historica}
\addcontentsline{toc}{chapter}{Historical Note}
\markboth{Historical Note}{}

\thispagestyle{empty}

\noindent This book is a work of fiction. Joaquim Antônio Correia, Patient
1,260, and the creature he carries never existed. The institution narrated
as \textit{the Colony} is fictional.

\vspace{0.8em}

\noindent But the ground upon which this fiction was built is real --- and
that is why it must end here, and not with a jump scare.

\vspace{0.8em}

\noindent During the 1915 drought, the government of Ceará created the
\textbf{Alagadiço Camp}, near Fortaleza, where retirantes were fenced in
and confined. The \textbf{Patu Camp}, in Senador Pompeu, belongs to a
different historical moment: the 1932 drought, with operation situated in
1932--1933. This edition adopts Alagadiço as the referent for its 1915
setting and retains Joaquim, his family, the trench, and the 1917 transfer
as fictional elements. Totals of internees and deaths are not presented as
precise.

\vspace{0.8em}

\noindent The \textbf{Barbacena Hospital-Colony}, in Minas Gerais, truly
existed. Over the course of the twentieth century, an estimated
\textbf{60,000 people} died within its walls --- many of them with no
psychiatric diagnosis whatsoever: they were poor, Black, epileptic,
single mothers, homosexuals, ``inconvenient'' people. There was
electroshock without anesthesia, restraint, hunger, mass graves, and the
sale of bodies to medical schools. None of this is invented by this book;
it is documented history --- see, among other sources, \textit{Brazilian
Holocaust}, by Daniela Arbex.

\vspace{0.8em}

\noindent The supernatural horror of these pages is a metaphor. The pain
that crosses generations, that no one processes and that passes on like a
contagion, is the only true creature of the book --- and it has a name:
forgetting. The worst thing one can do with the dead of the pens and the
Colony is not to haunt them with a story. It is to leave them with none.

\vspace{0.8em}

\noindent \textbf{On the risk of aestheticization:} the authors are aware
that trauma literature walks the frontier between symbolic elaboration and
the consumption of others' suffering. This book does not intend to offer
the horror of the Colony as spectacle, nor to replace the testimony of the
real victims and their descendants --- which remains irreplaceable and
prior to any fiction. If reading these pages produces only scares, it will
have failed; if it produces discomfort in the face of forgetting, it will
have fulfilled its function. To the reader who feels affected: interrupting
the reading is always legitimate, and support is available --- CVV 188
(24 hours, Brazil) and the 988 Suicide \& Crisis Lifeline (United States).

\vspace{1.2em}

\noindent\hfill\textit{So that what time buried may not be erased.}\hspace{2em}
